% Phase 2 limitations (Hardik).
\section{Limitations}
\label{sec:limitations}
The Phase 1 limitations carry over unchanged: Tox21 assays are in-vitro screens, not
organism toxicity; labels are noisy (44 conflicting duplicate pairs) and heavily missing
(17\% of the matrix, and not missing at random); several endpoints have very few positives;
and a scaffold split, while much harder than a random split, is still not prospective
validation. On top of these, the Phase 2 additions have their own caveats:

\begin{itemize}
\item \textbf{One tuning seed.} Hyper-parameters are chosen with seed 0 on
\texttt{scaffold\_det} validation only; differences between neighbouring configurations
that are smaller than the seed-to-seed spread reported in the tables should not be
over-interpreted.
\item \textbf{Hybrid parents share their history.} The fusion and embedding hybrids
warm-start from the MPNN checkpoint of the same split and seed, which itself was
early-stopped on that split's validation. This is the same information the GNN parent
had, so the comparison against the parent is fair, but the encoder has seen validation
through early stopping and so have all models in this comparison.
\item \textbf{Ablations are MPNN-based.} The fusion graph branch and the embedding
features use the tuned MPNN only; conclusions about combining branches may not transfer
to GCN or AttentiveFP encoders.
\item \textbf{Statistics are exploratory.} The Wilcoxon tests run over 12 endpoints and
many model pairs without multiplicity correction; $p$ values near 0.05 are suggestions.
\item \textbf{Compute budget.} All training is single-machine CPU; the GNN hyper-parameter
grid (36 configs) and the hybrid runs are small by deep-learning standards.
\end{itemize}
